\section*{Capítulo \ref{ch:g7:volcanoes} --- Os vulcões}

\begin{solution}{exo:g7:volcanoes:1}
(a) \omterm{def:g7:volcanoes:vent}{cratera}; (b) \omterm{def:g7:volcanoes:magma}{magma}; (c) \omterm{def:g7:volcanoes:vent}{boca}; (d) \omterm{def:g4:volcanoes-earthquakes-first-look:eruption}{lava}; (e) \omterm{def:g7:volcanoes:magma-chamber}{câmara magmática}.
\end{solution}

\begin{solution}{exo:g7:volcanoes:2}
(a) Efusiva. (b) Explosiva. (c) Efusiva: uma fonte de lava lança gotas
incandescentes que caem de volta e correm; o \omterm{def:g7:volcanoes:magma}{magma} não é transformado em
cinzas. (d) Explosiva: é um \omterm{def:g7:volcanoes:pyroclastic-flow}{fluxo piroclástico}.
\end{solution}

\begin{solution}{exo:g7:volcanoes:3}
(a) Verdadeiro. (b) Falso: o gás vulcânico mais abundante é o vapor de
água; o gás carbônico vem depois. (c) Falso: um \omterm{def:g7:volcanoes:pyroclastic-flow}{fluxo piroclástico} é uma
avalanche de gás quente, de cinzas e de pedaços de \omterm{def:g3:stones-pebbles-sand:rock}{rocha}. (d) Falso: eles
ficam dentro da \omterm{def:g7:plate-tectonics:plate}{placa} do Pacífico, a mais de 3200~km do limite mais
próximo, sobre um \omterm{def:g7:volcanoes:hotspot}{ponto quente}.
\end{solution}

\begin{solution}{exo:g7:volcanoes:4}
(a) De cerca de 1000 a 1200~graus. (b) De cerca de 600 a 800~graus.
(c) Perto da \omterm{def:g7:volcanoes:vent}{boca}, onde a \omterm{def:g4:volcanoes-earthquakes-first-look:eruption}{lava} é amarela (de 1000 a 1200~graus), e não
laranja (de 800 a 1000~graus): a \omterm{def:g4:volcanoes-earthquakes-first-look:eruption}{lava} esfria à medida que se afasta da
\omterm{def:g7:volcanoes:vent}{boca}.
\end{solution}

\begin{solution}{exo:g7:volcanoes:5}
(a) $838 + 221 = 1059$ \omterm{def:g4:volcanoes-earthquakes-first-look:volcano}{vulcões}. (b) $1059 \div 1214 \approx 0.872$: cerca
de \textbf{87~\%}. (c) $69 + 18 + 12 + 0.5 = 99.5$, cerca de 100~\%: a
pequena diferença vem do arredondamento de cada porcentagem.
\end{solution}

\begin{solution}{exo:g7:volcanoes:6}
(a) $564 \div 660 \approx 0.855$: cerca de \textbf{85~\%}. (b)
$660 - 564 - 62 - 33 = 1$: um cone íngreme. (c) Os cones íngremes são
construídos por \omterm{def:g4:volcanoes-earthquakes-first-look:eruption}{lava} viscosa e cinzas, a marca das \omterm{def:g7:volcanoes:eruption-types}{erupções explosivas}: ao
longo do Círculo de Fogo, onde uma \omterm{def:g7:plate-tectonics:plate}{placa} afunda embaixo de outra,
esperamos muitas \omterm{def:g7:volcanoes:eruption-types}{erupções explosivas}, como as do monte Santa Helena e do
Pinatubo.
\end{solution}

\begin{solution}{exo:g7:volcanoes:7}
(a) $1170 - 1140 = 30$~graus. (b) $30 \div 10 = 3$~graus por quilômetro.
(c) No túnel, a \omterm{def:g4:volcanoes-earthquakes-first-look:eruption}{lava} não toca o ar: o teto endurecido acima dela segura o
calor, como a tampa de uma panela. Ao ar livre, a sua superfície perde
calor para o ar na hora.
\end{solution}

\begin{solution}{exo:g7:volcanoes:8}
(a) Kauai, por pelo menos $5.5 - 0.7 = 4.8$~milhões de anos. (b) Kauai: a
\omterm{def:g4:clouds-rain-snow:precipitation}{chuva} e as ondas a desgastam há muito mais tempo. (c) A Ilha Grande, onde
ainda sai \omterm{def:g4:volcanoes-earthquakes-first-look:eruption}{lava} nova. (d) Para o noroeste: Kauai, levada para longe do
\omterm{def:g7:volcanoes:hotspot}{ponto quente} primeiro, hoje fica a noroeste dele.
\end{solution}

\begin{solution}{exo:g7:volcanoes:9}
(a) Quase todos os \omterm{def:g4:volcanoes-earthquakes-first-look:volcano}{vulcões} ficam em \omterm{def:g7:plate-tectonics:plate-boundary}{limites de placas} (87~\% na lista):
\omterm{def:g4:volcanoes-earthquakes-first-look:volcano}{vulcões} tão longe de qualquer limite são uma exceção. (b) Bem fundo no
\omterm{def:g6:inside-the-earth:layers}{manto}, embaixo do Havaí, um \omterm{def:g7:volcanoes:hotspot}{ponto quente} continua quente durante milhões
de anos. Ele derrete em parte a \omterm{def:g3:stones-pebbles-sand:rock}{rocha} acima dele. O \omterm{def:g7:volcanoes:magma}{magma}, mais leve que a
\omterm{def:g3:stones-pebbles-sand:rock}{rocha} sólida em volta, sobe através da \omterm{def:g7:plate-tectonics:plate}{placa} e entra em \omterm{def:g4:volcanoes-earthquakes-first-look:eruption}{erupção}. Enquanto
a \omterm{def:g7:plate-tectonics:plate}{placa} se move, novos \omterm{def:g4:volcanoes-earthquakes-first-look:volcano}{vulcões} crescem sobre o \omterm{def:g7:volcanoes:hotspot}{ponto quente} e os velhos
são levados embora. (c) Pontos laranja (``dentro de uma \omterm{def:g7:plate-tectonics:plate}{placa}'').
\end{solution}

\begin{solution}{exo:g7:volcanoes:10}
O primeiro \omterm{def:g7:volcanoes:magma}{magma} tinha perdido a maior parte do seu gás no caminho para
cima, como uma garrafa de água com gás deixada aberta ou parada: quando ele
chegou à superfície, não sobrava nada para despedaçá-lo, e ele saiu
devagar. O \omterm{def:g7:volcanoes:magma}{magma} de 12 e 15~de~junho ainda guardava o seu gás, como a
garrafa sacudida: quando a pressão caiu perto da superfície, o gás formou
bolhas de uma vez, que incharam e rasgaram o \omterm{def:g7:volcanoes:magma}{magma} viscoso em cinzas. O
primeiro \omterm{def:g7:volcanoes:magma}{magma} é o que tinha perdido o seu gás.
\end{solution}

\begin{solution}{exo:g7:volcanoes:11}
(a) $10 \div 80 = 0.125$~hora, e $0.125 \times 60 = 7.5$~minutos. (b)
$10 \div 130 \approx 0.077$~hora, cerca de $0.077 \times 60 \approx
4.6$~minutos, então cerca de 5~minutos. (c) A família precisa de
20~minutos, mas o fluxo chega em 5 a 8~minutos: quando o fluxo é visto, já
é tarde demais. A família só pode estar segura se partir quando os
cientistas virem os sinais de aviso, antes da \omterm{def:g4:volcanoes-earthquakes-first-look:eruption}{erupção}.
\end{solution}

\begin{solution}{exo:g7:volcanoes:12}
(a) $111 \times 360 = 39\,960$: cerca de \textbf{40\,000~km}. (b)
$40\,000 \div 15 \approx 2667$: cerca de 2\,670~km por dia. (c)
$2667 \div 24 \approx 111$: cerca de \textbf{110~km por hora}. (d) Pelo
vento: as cinzas finas não têm velocidade própria depois que a \omterm{def:g4:volcanoes-earthquakes-first-look:eruption}{erupção} as
lançou para cima; elas flutuam no alto do ar e o vento as leva, enquanto
elas caem devagar.
\end{solution}

\begin{solution}{pb:g7:volcanoes:1}
\textbf{1.} De 16 a 31~de~março: $31 - 16 = 15$~dias; depois, 30~dias de
abril e 18~dias de maio: $15 + 30 + 18 = 63$~dias.

\textbf{2.} $10\,000 \div 62 \approx 161$: mais de 160 pequenos \omterm{def:g4:volcanoes-earthquakes-first-look:earthquake}{terremotos}
por dia, em média.

\textbf{3.} O inchaço mostrava que o \omterm{def:g7:volcanoes:magma}{magma} tinha subido bem alto dentro da
\omterm{def:g1:landscapes-around-us:landform}{montanha} e empurrava a \omterm{def:g3:stones-pebbles-sand:rock}{rocha} para fora. Ele também deixava a face norte
íngreme e instável, pronta para deslizar.

\textbf{4.} Um triângulo azul fica onde uma \omterm{def:g7:plate-tectonics:plate}{placa} afunda embaixo de outra.
Esperamos um \omterm{def:g7:volcanoes:magma}{magma} viscoso e rico em gás, \omterm{def:g7:volcanoes:eruption-types}{erupções explosivas} e um cone
íngreme (um estratovulcão) --- exatamente o que mostra a primeira
fotografia.

\textbf{5.} $240 \div 60 = 4$~quilômetros por minuto.

\textbf{6.} A face inchada segurava o \omterm{def:g7:volcanoes:magma}{magma}, sob pressão, como a tampa de
uma garrafa. Quando ela deslizou, a pressão caiu na hora; o gás dissolvido
no \omterm{def:g7:volcanoes:magma}{magma} virou bolhas de uma vez e jorrou pela face aberta --- de lado,
onde a tampa tinha saído.

\textbf{7.} $27 \div 480 \approx 0.056$~hora, e $0.056 \times 60 \approx
3.4$~minutos. Como a explosão andou pelo menos tão rápido, ela levou no
máximo cerca de 3.4~minutos.

\textbf{8.} $8 \div 130 \approx 0.062$~hora, e $0.062 \times 60 \approx
3.7$: cerca de \textbf{4~minutos}.

\textbf{9.} $24 \div 8.8 \approx 2.7$: a coluna subiu cerca de 2.7 vezes
mais alto que o monte Everest.

\textbf{10.} Efusivas: a \omterm{def:g4:volcanoes-earthquakes-first-look:eruption}{lava} viscosa saiu sem explodir e se empilhou sobre
a \omterm{def:g7:volcanoes:vent}{boca}. Como a garrafa que ficou parada, esse \omterm{def:g7:volcanoes:magma}{magma} já tinha perdido a
maior parte do seu gás.

\textbf{11.} Antes, a \omterm{def:g1:landscapes-around-us:landform}{montanha} era um cone liso e pontudo coberto de \omterm{def:g4:clouds-rain-snow:precipitation}{neve},
acima de florestas verdes. Depois, o seu cume e a sua face norte
desapareceram: uma \omterm{def:g7:volcanoes:vent}{cratera} enorme aberta na nossa direção tomou o seu
lugar, e as encostas e o \omterm{def:g1:landscapes-around-us:landform}{vale} na frente estão cinzentos de cinzas e de
\omterm{def:g3:stones-pebbles-sand:rock}{rocha} quebrada.

\textbf{12.} $2950 - 2549 = 401$: a \omterm{def:g1:landscapes-around-us:landform}{montanha} perdeu cerca de
\textbf{401~m} de altura, cerca de 400~metros, como foi contado no 4º~ano.
\end{solution}
